\chapter{Background}
\label{ch:background}

This chapter reviews the background needed to treat \emph{refactoring as a process}, rather than only as a change between two endpoints. It begins with definitions of refactoring and refactoring sequences (\S\ref{sec:def-framing}), then discusses outcome-based evaluation and the code-quality signals commonly used to judge whether code has improved (\S\ref{sec:outcome-eval}--\S\ref{sec:smells-metrics}). It then reviews refactoring mining, recommendation and sequence synthesis, and IDE-based process tracing (\S\ref{sec:refactoring-mining}--\S\ref{sec:process-aware-logging}). The chapter ends by positioning this project against that work (\S\ref{sec:background-summary}).

\section{Definitions and framing}
\label{sec:def-framing}

\subsection{Refactoring as behaviour-preserving transformation}
Refactoring is usually defined as changing a program's internal structure to improve qualities such as readability or maintainability, while preserving externally observable behaviour. Classic definitions stress that refactoring is built from small, semantics-preserving steps, with safety supported by preconditions and validation such as testing \cite{opdyke1992refactoring,fowler1999refactoring,fowler2018refactoring}. Opdyke's work frames refactorings as program transformations with applicability conditions intended to preserve behaviour under a static approximation \cite{opdyke1992refactoring}. Fowler's catalogue helped popularise refactoring as an incremental developer activity made up of many small operations \cite{fowler1999refactoring,fowler2018refactoring}. Mens and Tourw\'e later survey the wider research landscape, bringing together definitions, taxonomies, and tool support, and emphasising the role of precondition checking and program analysis in automated refactoring \cite{mens2004survey}.

In practice, this clean definition is harder to apply. Developers often refactor while also adding features or fixing bugs, and some refactorings are only partly completed. As a result, it is often ambiguous whether a change seen in version history is a ``pure refactoring''. In this thesis, behaviour preservation is therefore treated as an \emph{intended property} of refactoring steps, while recognising that real traces may include edits with more than one purpose.

\subsection{Atomic refactorings, smells, and sequences}
Research on refactoring often describes refactorings as a set of \emph{operators}, or atomic refactorings, such as \textsc{Extract Method}, \textsc{Move Method}, and \textsc{Rename} \cite{fowler1999refactoring,mens2004survey}. Larger changes can then be understood as sequences of these smaller operations. Code smells are closely related: they are heuristic indicators of possible design problems, and therefore of possible refactoring opportunities \cite{fowler1999refactoring,fowler2018refactoring}. Although smells are not formal correctness criteria, they provide a useful bridge between developer judgement and measurable signals, discussed further in \S\ref{sec:smells-metrics}.

A key distinction in this thesis is between judging only the start and end states of a refactoring, and judging the path taken between them. Two developers may reach the same final code, but one path may involve more temporary degradation, churn, or disruption than the other. These differences can matter for productivity and risk, even when the endpoint looks the same.

\subsection{Representing refactoring as states and actions}
To represent this refactoring path explicitly, we use a state/action/trace model. A codebase snapshot is treated as a \emph{state} $s \in \mathcal{S}$, such as the repository at a particular point in time. An \emph{action} $a \in \mathcal{A}$ is a developer edit step, which may be an IDE-supported refactoring, a manual change, or a mixture of both. A refactoring process is then a \emph{trace}
\[
\tau = (s_0, a_0, s_1, a_1, \dots, a_{T-1}, s_T),
\]
where each action $a_t$ transforms state $s_t$ into $s_{t+1}$. This makes it possible to talk about ``process quality'' as a property of the whole trace $\tau$, rather than only as a comparison between the first and final states $(s_0, s_T)$.

This framing is also close to recent software-engineering work that models agent and developer behaviour as a partially-observable Markov decision process, where a trajectory is a sequence of action and observation pairs \cite{gautam2025refactorbench}. We keep the explicit state notation because this thesis scores the intermediate code states themselves, not only the actions or observations produced along the way.

\section{Outcome-based refactoring evaluation}
\label{sec:outcome-eval}

\subsection{Judging refactoring by the final state}
Refactoring is often evaluated by comparing quality indicators before and after the change. Under this view, a refactoring is successful if the final snapshot improves on the starting snapshot according to selected measures, such as lower coupling, higher cohesion, lower complexity, or fewer code smells. This endpoint-focused view appears both in empirical studies of refactoring impact and in recommendation systems that rank candidate refactorings by their predicted improvement.

A clear example is \emph{search-based refactoring}, where refactoring is treated as an optimisation problem over possible refactoring operations or sequences. Harman and Tratt apply multi-objective search at the design level, treating coupling and cohesion as objectives that may need to be traded off against each other \cite{harman2007pareto}. Later work extends this idea to many-objective settings and adds constraints or preferences for practical concerns such as test coverage or prioritisation \cite{mohan2019manyobjective}. This line of work is useful here because it shows that ``better'' code is not a single fixed property: different quality goals can pull in different directions.

The same endpoint-based idea is also used in refactoring recommendation systems. These systems compare possible refactorings and rank them by the improvement they are expected to produce. Some go beyond code-structure metrics by using information such as requirements traceability, since a design may be easier to maintain when it also supports tasks like impact analysis \cite{wang2018traceability}. More recent work has also looked at how to choose among large sets of candidate refactorings, and which code characteristics are associated with positive or negative effects in practice \cite{nikolaidis2024metricsbased}.

\subsection{What endpoint evaluation captures-and what it misses}
Endpoint-based evaluation is useful when the main concern is whether the final code is better than the starting code under some metric $Q$. It gives recommendation systems a simple strategy: simulate each candidate refactoring and choose the one that gives the largest improvement in $Q$. It is also useful for empirical studies, where refactoring events can be compared with changes in measured code quality. This fits naturally with many existing metrics, since they are defined over a single codebase snapshot.

The weakness of this approach is that it says little about the refactoring process itself. Two developers may reach the same endpoint through very different paths, but endpoint metrics would treat those paths as equivalent. In practice, the intermediate states can still matter: code may become harder to understand for a period of time, the modular structure may be disrupted, or the build and tests may break before being repaired. The path can also be costly in other ways, for example by creating extra churn, making review harder, or increasing integration risk. Empirical work on modularity improvement shows this trade-off clearly: improvements in design metrics can come with substantial structural disruption \cite{paixao2017balancing}. This motivates evaluating not only the final code, but also the stability and cost of the path taken to reach it.

This is the point at which the process view becomes useful. Final-code metrics still matter, but they need to be considered alongside what happened on the way there: whether the work introduced temporary degradation, required rework, created unnecessary churn, or broke build/test stability.

\section{Code smell detection and quality metrics}
\label{sec:smells-metrics}

\subsection{Using code smells as design signals}
Code smells provide a way to connect qualitative design concerns with measurable indicators \cite{fowler1999refactoring}. Smell detectors usually identify these cases either through hand-written rules or learned models. They often rely on structural properties such as size, coupling, and cohesion, with some tools also using lexical or semantic features. One prominent rule-based example is DECOR, which defines smells and design anti-patterns through measurable symptoms and detection rules \cite{moha2010decor}. DECOR is widely used in empirical studies because its definitions are explicit and it can be run efficiently at scale.

Smell detectors are useful, but they should not be treated as ground truth. Comparative studies show that tools can disagree in both precision and recall, and that performance varies across systems and smell types \cite{paiva2017evaluation}. Much of this comes from the practical choices each tool makes, such as thresholds, feature sets, and implementation details. Detectors can also produce false positives, marking code as ``smelly'' even when there is no clear negative impact or when the design reflects an acceptable trade-off \cite{arcellifontana2016falsepositives}. For this reason, smell counts are used in this thesis as one signal among several. They are combined with structural metrics, readability-oriented measures, and process-level safety signals, rather than being allowed to determine the process-quality score on their own.

\subsection{Structural and readability metrics}
Static metrics provide another way to measure code quality at the level of classes, methods, and packages. Common metric families cover coupling, cohesion, complexity, and size, and are often used as indicators of maintainability. These structural measures are useful, but they do not capture everything developers care about when refactoring. Readability is one example: models that combine structural features, such as line structure, with textual features, such as identifier patterns, tend to align better with human judgements of readability \cite{scalabrino2018readability}. This matters for process evaluation because readability and understandability are common reasons for refactoring, and they can change noticeably in intermediate states even if the final code improves.

\subsection{What this means for the process-quality metric}
The process-quality metric $J(\tau)$ uses these signals, but it also has to account for their limitations. The literature suggests three design choices:
\begin{enumerate}
  \item \textbf{Use multiple signals.} No single metric or smell detector gives reliable ground truth across all contexts \cite{paiva2017evaluation,arcellifontana2016falsepositives}. Combining several signals reduces dependence on the quirks of one tool.
  \item \textbf{Separate endpoint improvement from process cost.} Improving the final snapshot is important, but disruption and temporary degradation can still matter even when the endpoint improves \cite{paixao2017balancing}.
  \item \textbf{Keep the components interpretable.} A process critique is more useful when it can point to a specific issue, such as a readability dip, churn spike, or smell increase, rather than only reporting a single opaque score.
\end{enumerate}

A process metric should also include some notion of developer \textbf{effort}, such as the number of steps in the trajectory or the size of the edits between successive states. Fewer steps are not always better, since small steps can reduce risk, but step count and edit size are still useful for identifying unnecessary detours and rework. Similar cost terms are already used alongside structural quality goals in search-based refactoring and recommendation work, for example when trying to improve quality while minimising change \cite{harman2007pareto,mohan2019manyobjective}.

Table~\ref{tab:metric-candidates} summarises the candidate signals, their main benefits and limitations, and examples of Java tooling. The table does not define the final formula for $J(\tau)$, including weights or normalisation. Its purpose is to motivate the signals used later; the precise definition of $J(\tau)$ is given in the methodology chapter and used consistently in the evaluation.

\begin{table}[htbp]
\centering
\small
\begin{tabular}{L{0.19\textwidth} L{0.26\textwidth} L{0.28\textwidth} L{0.19\textwidth}}
\hline
\textbf{Signal} & \textbf{Pros} & \textbf{Cons / threats} & \textbf{Typical tooling} \\
\hline
Structural metrics (coupling/cohesion/\newline complexity/size) &
Widely used; efficient; interpretable at class/method granularity &
Validity depends on context; may incentivise disruptive restructuring; metric interactions &
Static analysis metric extractors; IDE/CI integrations \\
\hline
Smell counts / severity (rule-based) &
Directly targets design issues; aligns with developer's motivations &
Detector disagreement and false positives; threshold sensitivity; context dependence &
DECOR-style detectors \cite{moha2010decor}; other smell tools \cite{paiva2017evaluation} \\
\hline
Readability models &
Closer to human comprehension than structural-only metrics; captures lexical/style clues &
Model generalisation and dataset bias; may be language/style dependent &
Readability predictors \cite{scalabrino2018readability} \\
\hline
Churn / disruption (LOC changed, file moves/renames) &
Captures cost and instability of change; naturally trajectory-oriented &
Churn may be necessary for large refactors; can penalise required restructuring &
VCS mining; diff statistics \\
\hline
Process length / effort (number of steps; edit magnitude) &
Trajectory-native cost signal; discourages detours/rework; easy to compute &
Sensitive to step granularity (commits vs.\ checkpoints vs.\ IDE events); may penalise safe micro-steps; requires calibration &
VCS checkpoints; IDE logs; diff/AST edit size \\
\hline
Build/test preservation signals &
Direct safety signal; aligns with ``behaviour-preserving intent'' &
Requires runnable builds/tests; noisy due to flaky tests; environment dependence &
CI logs; local build/test runs \\
\hline
Traceability-augmented signals &
Reflects maintenance tasks beyond code structure; complements metric-only ranking &
Requires trace links (often unavailable); domain-specific assumptions &
Traceability-based ranking methods \cite{wang2018traceability} \\
\hline
\end{tabular}
\caption{Candidate signals for constructing a process-quality objective. We combine endpoint quality change with trajectory-oriented cost and stability terms, rather than relying on any single indicator.}
\label{tab:metric-candidates}
\end{table}


\section{Refactoring detection and mining from code changes}
\label{sec:refactoring-mining}

\subsection{From diffs to refactoring actions}
To evaluate a process trace in practice, developer actions either need to be captured directly, for example through IDE event logs, or inferred from the differences between snapshots such as bursts of manual user edits. Much of this work uses the second route: given two versions of a program, infer whether a refactoring happened and identify which refactoring types occurred.

\emph{Structured code differencing} is important here because many refactorings are not easy to recognise from line-based diffs alone. AST-based differencing can represent moves, renames, and other non-local edits more directly. GumTree, for example, maps fine-grained AST nodes between program versions and produces edit scripts that better reflect the structural changes developers intended \cite{falleri2024gumtree}. This kind of differencing is often used as the base layer for refactoring detection.

Modern refactoring miners use these structural mappings together with heuristics for particular refactoring types. RefactoringMiner detects a wide range of refactorings between two versions of the code by matching related classes, methods, and fields \cite{alikhanifard2025rminer}. RefDiff takes a similar approach, using similarity heuristics and static analysis to detect common refactoring types in version histories \cite{silva2017refdiff}. For process analysis, tools like these are useful because they can turn a transition $s_t \rightarrow s_{t+1}$ into a set of detected refactoring actions, giving the trace a concrete \emph{action vocabulary}.

\subsection{What refactoring detectors miss}
Refactoring detection tools usually report what kind of refactoring occurred, such as \textsc{Move Method} or \textsc{Extract Class}, and which classes, methods, or fields were involved. This fits the operator-based view of refactoring used in refactoring guides and surveys \cite{fowler1999refactoring,mens2004survey}. However, the literature also reports several limitations that matter for process evaluation.

\paragraph{Mixed and tangled changes.}
Version control commits often mix refactoring with feature work or bug fixes, and they may bundle several unrelated activities together. A single commit-level transition can therefore include both behaviour-preserving and behaviour-changing edits. In these cases, refactoring detection may return only partial or ambiguous results. This makes snapshot-derived actions harder to interpret and motivates careful trace construction, such as using finer checkpoints and methods that can tolerate imperfect refactoring detections.

\paragraph{Limited coverage and combined refactorings.}
Detectors only cover a finite set of operators, and refactorings can be composed, such as a move followed by a rename, in ways that make matching harder. Some refactorings are also performed manually without tool support and may not fit the patterns that detectors expect. For process evaluation, this means that ``actions'' should be treated as partially observed: detected refactorings are useful features, but they are not a complete description of developer intent.

\paragraph{Behaviour preservation is not verified.}
Detection is usually structural. It infers that a change \emph{looks like} a refactoring, but it does not prove that behaviour was preserved. Because of this, downstream evaluation should not treat a detected refactoring as automatically safe. Where possible, it should also use independent safety signals, such as whether the build and tests still pass.


\subsection{Tie-in to our work}
Refactoring mining is important infrastructure for building traces and describing actions at scale \cite{alikhanifard2025rminer,silva2017refdiff}. This thesis does not try to improve refactoring detection accuracy. Instead, mined refactorings are used as part of the observational layer: they provide labels and features for transitions between states, which can then support process scoring and reference-trajectory generation. The next section builds on this by discussing work that generates and navigates refactoring sequences.



% =========================
% Background (continued)
% Sections 2.5--2.7
% =========================

\section{Refactoring recommendation and search/synthesis}
\label{sec:recommendation-synthesis}

The previous sections point to two needs for process-level evaluation: a way to represent and score a developer trace, and a way to construct alternative trajectories for comparison. There is already a large body of work on generating refactoring suggestions and sequences, but most of it is not designed to evaluate a developer's process. This section reviews three related families of work: synthesis approaches that reconstruct or complete refactoring sequences, search-based approaches that optimise sequences under quality objectives, and interactive approaches that adapt recommendations based on developer feedback.

\subsection{Synthesis and path completion from partial edits}
One way to generate alternatives is to treat refactoring as a program transformation problem. Given a start program and a target, or a partially edited program, the task is to synthesise a sequence of refactoring operators that carries out the transformation. ReSynth, for example, synthesises a sequence of IDE-supported refactorings that is consistent with a developer's partial edits \cite{raychev2013resynth}. It uses the edited classes, methods, or fields, together with the operators' preconditions, to constrain the search. Its goal is reachability and consistency: can some operator sequence reproduce the developer's edits? This is different from evaluating process quality, but it shows that refactoring can be treated as a structured search problem and that the search can be guided by evidence from the developer trace.

A related line of work learns transformations from examples. REFAZER learns program transformations from input-output examples of code edits \cite{rolim2017refazer}. Rather than assuming a fixed list of defined refactoring types, it synthesises AST-level rewrite rules from a small number of examples and can apply them to automate repetitive edits. Although REFAZER is not a refactoring miner in the narrow Fowler sense, it is relevant here because it shows that plausible transformations can be inferred from the edits themselves, including edits that do not fit neatly into IDE-supported refactoring operations.

\paragraph{Relevance and limitations for process evaluation.}
Synthesis systems are mainly concerned with whether a target program can be reached, or whether a sequence can be found that is consistent with partial edits, while still satisfying refactoring preconditions \cite{raychev2013resynth}. They do not judge the intermediate choices the developer actually made: for example, whether those choices introduced avoidable churn, temporary degradation, or extra risk. They are still useful for this thesis because they show how candidate alternative trajectories can be generated in a structured way, rather than invented by hand.

\subsection{Search-based optimisation of refactoring sequences}
A second family treats refactoring recommendation as an optimisation problem. Search-based methods explore possible changes or sequences of changes, then choose those that perform best against one or more quality objectives \cite{harman2007pareto}. Harman and Tratt, for example, optimise for multiple objectives at once, treating coupling and cohesion as competing concerns and producing Pareto fronts rather than reducing design quality to a single score \cite{harman2007pareto}. This is useful because it makes trade-offs between different notions of ``better code'' explicit.

Later work brings in concerns that are closer to everyday development. Many-objective approaches, for example, consider not only quality improvement, but also prioritisation and how refactoring effort is spread across the codebase \cite{mohan2019manyobjective}. This matters for process evaluation because it recognises that refactoring has costs while it is being carried out, not only benefits in the final version of the code.

Search-based techniques also show that improving metrics can come with disruption. Empirical studies of modularity optimisation show that maximising cohesion and coupling improvements can substantially disturb the original modular structure \cite{paixao2017balancing}. More recent work makes this trade-off explicit by treating \emph{architectural distance}, a measure of disruption, as an objective alongside structural quality \cite{cortellessa2023manyobjective}. Other work adds socio-technical costs, such as reviewer availability, to the recommendation process \cite{chen2024morcora}. This matters here because a final design can look better under structural metrics while still being costly or disruptive to reach.

\paragraph{Relevance and limitations for process evaluation.}
Search-based refactoring is useful here because it provides two ideas this thesis builds on: explicit quality objectives and generated candidate sequences. The difference is in how those sequences are used. In most recommendation work, the generated sequence is the proposed solution, and success is measured by the improvement it produces or by whether developers prefer it. In this thesis, the developer's observed trajectory is what we evaluate. Generated alternatives are used only as references for comparison.

\subsection{Goal-directed navigation and architectural refactoring}
A third family focuses on guiding refactoring toward a target design. Refactoring Navigator, for example, takes a current system and a desired architecture or design, then recommends a sequence of refactoring steps using a search-based strategy \cite{lin2016navigator}. The important point is that it guides developers step by step, rather than only ranking isolated refactorings. Controlled evaluations report large reductions in task time and manual edits compared with unguided refactoring \cite{lin2016navigator}, suggesting that sequence-level guidance can change how developers refactor in practice.

Goal-directed refactoring is relevant because it shows how an overall design goal can be translated into a sequence of concrete refactoring actions. This supports the comparison used in this thesis: an observed developer trajectory can be compared with a higher-scoring reference trajectory that works towards the same end state.

\paragraph{Relevance and limitations for process evaluation.}
Navigation systems usually focus on reaching a target design, using design metrics to guide the suggested steps \cite{lin2016navigator}. Their purpose is not to judge the quality of a developer's actual path, or to explain where that path moved away from a better alternative. Even so, they are important for this thesis because they show that refactoring can be treated as a sequence of decisions, and that sequence-level guidance can be computed and shown to developers in a usable form.

\subsection{Interactive and developer-in-the-loop recommendation}
Purely automatic optimisation can end up misaligned with developer intent, constraints, or context. Interactive refactoring recommenders address this by incorporating user feedback during the search process. Alizadeh and Kessentini propose an interactive and dynamic search-based approach that presents refactoring recommendations iteratively and updates the search based on developer decisions \cite{alizadeh2020interactive}. This highlights that refactoring recommendation is not only about finding a high-scoring sequence under fixed metrics; it also involves keeping recommendations aligned with developer preferences and constraints as they evolve during a task.

\paragraph{Relevance and limitations for process evaluation.}
Interactive recommenders are not evaluation tools, but they are closely aligned with the instrumentation and delivery mechanism we envision. If process evaluation is delivered in an IDE, the feedback has to be understandable, minimally disruptive, and sensitive to developer context. Work on interactive recommendation provides ideas for presenting sequence-level suggestions and for building feedback loops into the developer's workflow \cite{alizadeh2020interactive}. The remaining open problem is to shift from ``recommend to optimise'' to ``compare and evaluate an observed trajectory relative to alternative, counterfactual trajectories''.

\subsection{Beyond structural metrics: alternative objectives for ranking refactorings}
A recurring theme in recommendation research is that structural metrics alone may not be enough to capture what makes a refactoring ``good''. One example is ranking refactoring candidates using requirements traceability alongside code metrics. Approaches that include traceability treat ``quality'' as partly determined by how well requirements map to code, and they are motivated by maintenance tasks like impact analysis and comprehension \cite{wang2018traceability}. This supports the idea that process-quality objectives should not be assumed to be purely structural: the relevant costs and benefits of refactoring can include non-structural properties and context-specific constraints.

Recent work also looks at developer motivations for refactoring, and tries to align recommendations with what developers are likely to prioritise. Studies that mine refactoring operations and repository history investigate product/process signals that correlate with refactoring activity and with developer intent, including proposals that use LLMs to infer motivations from commit artefacts \cite{whydevelopersrefactor2020,metricsmotivation2024}. While these works mostly focus on \emph{when} and \emph{why} refactoring happens, rather than \emph{how well} it is carried out as a process, they suggest that developer-aligned objectives may need to incorporate process signals and contextual information.

\subsection{Synthesis: what sequence-generation work enables, and what remains missing}
Sequence generation and recommendation systems establish that it is feasible to:
\begin{itemize}
  \item generate refactoring sequences that realise a transformation or reach a target (synthesis/navigation) \cite{raychev2013resynth,lin2016navigator},
  \item optimise sequences under explicit metric objectives (search-based refactoring) \cite{harman2007pareto,mohan2019manyobjective},
  \item incorporate developer feedback into iterative suggestion mechanisms (interactive recommendation) \cite{alizadeh2020interactive},
  \item and broaden ``quality'' beyond structural metrics where additional artefacts exist (e.g., traceability) \cite{wang2018traceability}.
\end{itemize}
However, these systems do not directly answer the two central questions posed by process-level evaluation:
\begin{enumerate}
  \item \emph{Was the developer's observed trajectory good under an explicit process-quality criterion?}
  \item \emph{Where did the developer's trajectory diverge from a higher-quality trajectory, and can that divergence be explained?}
\end{enumerate}
The distinction is subtle but important. Recommendation and synthesis work typically treats the algorithm's output as the primary object and evaluates its outcome or usefulness. Process evaluation instead treats the \emph{developer trace} as the primary object and seeks to compare it against counterfactual trajectories constructed under a process-quality objective.

\subsection{Recent trajectory-formalism work in software engineering}

Recent work has begun to formalise software-engineering agent behaviour as a trajectory of (action, observation) pairs in a POMDP-shaped environment \cite{gautam2025refactorbench}, with related work studying agent action patterns \cite{bouzenia2025trajectories} and training environments \cite{pan2025swegym}; a recent systematic literature review of LLM-based refactoring research surveys evaluation methodologies and open challenges in this area \cite{alomar2025slr}. These developments confirm that the trajectory framing is gaining traction in concurrent software-engineering research, but the work to date targets LM-agent behaviour and benchmark task success rather than human developer process quality. % TODO methodology cross-cite: gautam2025refactorbench failure-mode #2 (intermediate-broken-state cost) is direct support for the time-weighted W_BROKEN choice.
Empirical observational studies independently confirm that real-world refactoring proceeds as ordered sequences of operations interleaved with other change activities \cite{eilertsen2024refactorings}, motivating the trajectory-level framing we adopt.

\section{Process-aware work and IDE trace capture}
\label{sec:process-aware-logging}

Evaluating a refactoring process requires access to traces: sequences of intermediate states and actions. There are two broad sources of traces: (i) version control histories (commits/checkpoints), and (ii) IDE-level interaction logs that capture finer-grained behaviour. This section focuses on IDE logging and process-aware studies, since they show it is feasible to capture refactoring traces at a useful granularity, and they also highlight practical limitations that end up determining evaluation design.

\subsection{Logging developer interactions in IDEs}
Multiple studies have collected large-scale logs of developer interactions to mine workflows and behavioural patterns. Damevski et al.\ present an approach for mining sequences of developer interactions in Visual Studio using telemetry-like event streams, with the aim of discovering frequent patterns of usage without necessarily recording the full code content \cite{damevski2017usagesmells}. This is relevant for trace collection because it shows that useful traces can be captured as structured events (commands, navigation actions, edits) while still maintaining the developer's privacy. The flip side of the privacy-preserving design is that the trace cannot subsequently be replayed against the project state at each event, which limits what downstream analysis can do with it.

Poncin et al.\ combine heterogeneous software-repository logs (version control, bug trackers, mail archives) and apply process mining techniques to infer developer workflows \cite{poncin2011processmining}. This gives a clear example of turning low-level event records into higher-level activity models, such as sequences of commits, issue updates, and discussions. It also shows some of the practical challenges: event streams are noisy, developers multitask across repositories, and mapping raw events to semantically meaningful steps needs careful segmentation and interpretation. The events themselves sit at repository granularity rather than IDE granularity, so finer-grained signals such as which IDE command was run or whether a refactoring was invoked via the menu are not visible.

More recent infrastructure work has pivoted toward capturing developer-LLM interactions rather than refactoring-specific developer process traces. \emph{CodeWatcher} is a representative example: a lightweight client-server VS Code extension that captures fine-grained, semantically labelled IDE events (LLM-tool insertions, deletions, copy-paste, focus shifts) with timestamps to support post-hoc reconstruction of coding sessions \cite{basha2025codewatcher}. While the underlying telemetry mechanisms are directly applicable to refactoring traces, the 2023--2025 literature on IDE-side activity mining for refactoring specifically has been quiet-the broader process-mining-in-SE community has migrated toward repository and CI/CD logs rather than IDE event streams, and recent activity-recognition work targets LLM-coding patterns rather than refactoring sequences. The rich IDE-event capture that is being done targets VS Code and LLM coding rather than Java refactoring in IntelliJ, so the recordings do not transfer cleanly to refactoring-trajectory work. Classical IDE-side activity mining for refactoring trajectories therefore remains an area where canonical references are over seven years old.

Taken together, the available corpora either capture rich IDE events for the wrong activity, capture refactoring activity at repository rather than IDE granularity, or do not retain enough code content to replay the trajectory. Constructing a fit-for-purpose corpus is therefore a methodology question rather than a corpus-reuse question, and Chapter~\ref{ch:methodology} returns to it at \S\ref{sec:methodology-rater}.

\subsection{Granularity, segmentation, and semantic interpretation}
IDE logging raises the question of what counts as a ``step'' in a refactoring process. At commit granularity, steps are coarse and often tangled with non-refactoring changes. At event granularity, steps are extremely fine and need to be aggregated into coherent actions. The process mining and interaction analysis literature highlights that segmenting activity into tasks/episodes is not trivial \cite{poncin2011processmining,damevski2017usagesmells}. For refactoring in particular, a single conceptual refactoring can involve several micro-actions (navigate, edit, run a refactoring tool, adjust code, run tests). On the other hand, a short burst of edits might correspond to multiple actions.

A process-evaluation system therefore has to choose a step definition that is both (i) practical to extract reliably and (ii) meaningful enough to allow for interpretation.

\section{Summary and positioning: toward process-quality evaluation}
\label{sec:background-summary}

The literature discussed in \S\ref{sec:def-framing}--\S\ref{sec:process-aware-logging} establishes a strong foundation for modelling and measuring refactoring, but it also leaves open a clear space for process-level evaluation.

First, refactoring is well-defined as behaviour-preserving and supported by catalogues of operators and by tool-based precondition checking \cite{opdyke1992refactoring,fowler1999refactoring,mens2004survey}. Second, most evaluation work is still focused primarily on endpoints: it assesses refactoring success by looking at changes in quality indicators between start and end snapshots, including structural metrics and smell-based proxies \cite{harman2007pareto,moha2010decor}. These signals are useful, but they are imperfect and can disagree across tools, which motivates combining multiple signals and metrics \cite{paiva2017evaluation,arcellifontana2016falsepositives,scalabrino2018readability}. Third, refactoring mining and AST differencing provide practical ways to infer refactoring actions from code changes, which enables trace reconstruction from repository histories \cite{alikhanifard2025rminer,falleri2024gumtree}. Fourth, recommendation and synthesis systems show that it is feasible to generate refactoring sequences and plan toward final goal states \cite{raychev2013resynth,lin2016navigator,alizadeh2020interactive}, and that ``quality'' may need to include contextual objectives beyond just structural metrics \cite{wang2018traceability,whydevelopersrefactor2020,metricsmotivation2024}. Finally, IDE logging research shows that fine-grained traces can be captured and analysed, but it also highlights challenges around the segmentation and interpretation of them \cite{poncin2011processmining,damevski2017usagesmells,foutsekh2018interactiontraces}.

While there has been substantial research on \emph{detecting} refactorings, \emph{recommending} refactorings, and \emph{synthesising} refactoring sequences, as far as we are aware no prior system both (i) scores a developer's refactoring trace against an explicit process-quality metric and (ii) produces counterfactual reference trajectories tied to measurable divergence points. The closest existing work-sequence synthesis, search-based recommendation, and the recent trajectory-formalism studies surveyed above-generates or analyses sequences as the primary object, treating the algorithm's output as the solution rather than evaluating a developer's chosen path against alternatives. Existing research generally does not define and evaluate a criterion for ``good refactoring process'' that penalises intermediate degradation, disruption, or instability, nor does it provide a specific method to compare a developer's refactoring trajectory against alternative, higher-quality trajectories (relative to the chosen metric) and to explain divergence points from those trajectories.

We address this gap by treating refactoring as a trajectory optimisation and comparison problem: given a start state, an end state and an observed trace, defining a process-quality objective and constructing a higher-scoring reference trajectory for comparison. This leads to the overall thesis question: \emph{given the start and end state of a refactoring session, can we judge the quality of the path the developer took and point to places where a better path was available?}
